\chapter{Équations du Premier Degré \& Problèmes}
\label{chap:equations}

% ==============================================================================
% 1. OBJECTIFS & COMPÉTENCES DU RÉFÉRENTIEL
% ==============================================================================
\begin{cadreretenu}[Objectifs \& Compétences du DNB Pro]
\begin{itemize}
    \item \textbf{Connaissances} : Notion d'égalité, d'inconnue, de membre d'une équation, équation produit nul.
    \item \textbf{Capacités} :
    \begin{itemize}
        \item \capRealiser\ Résoudre algébriquement une équation de la forme $ax + b = c$ et $ax + b = cx + d$.
        \item \capRealiser\ Résoudre une équation produit nul de type $(ax + b)(cx + d) = 0$.
        \item \capAnalyser\ Mettre en équation un problème issu de la vie courante ou d'un domaine professionnel (choix de l'inconnue, mise en équation, résolution, conclusion).
        \item \capValider\ Tester et vérifier si un nombre est solution d'une équation.
    \end{itemize}
\end{itemize}
\end{cadreretenu}

\vspace{0.3cm}

% ==============================================================================
% 2. ACTIVITÉ D'INVESTIGATION
% ==============================================================================
\section{Activité d'Introduction : Le Principe de la Balance}

\begin{activite}{L'équilibre d'une balance à deux plateaux}
Une balance est en parfait équilibre. Sur le plateau de gauche, on a placé 3 masses identiques inconnues notées $x$ et une masse marquée de $5\text{ kg}$. Sur le plateau de droite, on a placé une masse $x$ et une masse marquée de $17\text{ kg}$.

\begin{center}
\begin{tikzpicture}[scale=0.85]
    % Socle et fléau
    \draw[fill=slate700] (-0.3,0) rectangle (0.3,2);
    \draw[fill=slate800] (-1.5,0) rectangle (1.5,0.3);
    \draw[line width=2pt, colPrimary] (-4,2) -- (4,2);
    \draw[fill=colSecondary] (0,2) circle (4pt);

    % Plateaux
    \draw[thick, slate600] (-4,2) -- (-4,1) -- (-2.5,1) -- (-4,2);
    \draw[thick, slate600] (4,2) -- (4,1) -- (2.5,1) -- (4,2);
    \draw[line width=1.5pt, slate800] (-4.5,1) -- (-2.2,1);
    \draw[line width=1.5pt, slate800] (2.2,1) -- (4.5,1);

    % Masses gauche
    \draw[fill=colExoBg, draw=colExo, thick] (-4.2,1) rectangle (-3.6,1.6) node[midway, font=\footnotesize\bfseries] {$x$};
    \draw[fill=colExoBg, draw=colExo, thick] (-3.5,1) rectangle (-2.9,1.6) node[midway, font=\footnotesize\bfseries] {$x$};
    \draw[fill=colExoBg, draw=colExo, thick] (-3.85,1.6) rectangle (-3.25,2.2) node[midway, font=\footnotesize\bfseries] {$x$};
    \draw[fill=amber!30, draw=amber!80!black, thick] (-2.8,1) rectangle (-2.3,1.5) node[midway, font=\tiny\bfseries] {$5\text{kg}$};

    % Masses droite
    \draw[fill=colExoBg, draw=colExo, thick] (2.5,1) rectangle (3.1,1.6) node[midway, font=\footnotesize\bfseries] {$x$};
    \draw[fill=amber!30, draw=amber!80!black, thick] (3.3,1) rectangle (4.2,1.9) node[midway, font=\footnotesize\bfseries] {$17\text{kg}$};
\end{tikzpicture}
\end{center}

\begin{enumerate}
    \item \capSApproprier\ Traduire l'équilibre de cette balance par une équation mathématique :
    \[ 3x + 5 = x + 17 \]
    \item \capAnalyser\ Si on retire une masse $x$ sur chaque plateau, que devient l'égalité ?
    \lignesreponse[2]
    \item \capRealiser\ Si on retire ensuite $5\text{ kg}$ sur chaque plateau, quelle masse pèsent les deux masses $x$ restantes ? En déduire la valeur de la masse $x$ :
    \lignesreponse[3]
\end{enumerate}
\end{activite}

% ==============================================================================
% 3. COURS ET NOTIONS ESSENTIELLES
% ==============================================================================
\section{Cours : Résolution Algébrique et Méthodes}

\subsection{Définitions et Propriétés Fondamentales}

\begin{cadredef}[Équation du premier degré à une inconnue]
Une \textbf{équation du premier degré} est une égalité comportant une lettre inconnue (souvent $x$) élevée à la puissance 1. \\
\textbf{Résoudre} une équation, c'est trouver toutes les valeurs de $x$ qui rendent l'égalité vraie. Chaque valeur trouvée est appelée une \textbf{solution}.
\end{cadredef}

\begin{cadreprop}[Règles de conservation de l'égalité]
Une égalité reste vraie si :
\begin{itemize}
    \item On \textbf{ajoute} ou on \textbf{soustrait} un même nombre aux deux membres :
    \[ A = B \ssi A + c = B + c \quad\text{et}\quad A - c = B - c \]
    \item On \textbf{multiplie} ou on \textbf{divise} les deux membres par un même nombre \textbf{non nul} ($c \neq 0$) :
    \[ A = B \ssi A \times c = B \times c \quad\text{et}\quad \frac{A}{c} = \frac{B}{c} \]
\end{itemize}
\end{cadreprop}

\subsection{Méthode Générale de Résolution}

\begin{cadremethode}[Résoudre une équation du type $ax + b = cx + d$]
\begin{enumerate}
    \item \textbf{Regrouper} tous les termes contenant $x$ dans le membre de gauche (en soustrayant $cx$).
    \item \textbf{Regrouper} toutes les constantes numériques dans le membre de droite (en soustrayant $b$).
    \item \textbf{Réduire} chaque membre pour obtenir une équation simplifiée de la forme $k x = m$.
    \item \textbf{Diviser} les deux membres par le coefficient $k$ ($k \neq 0$) : $x = \frac{m}{k}$.
    \item \textbf{Conclure} en vérifiant la solution trouvée.
\end{enumerate}
\end{cadremethode}

\begin{cadreexemple}[Résolution pas à pas]
Résoudre dans $\mathbb{R}$ l'équation $5x - 8 = 2x + 7$ :
\begin{align*}
    5x - 8 - 2x &= 7 && \text{(on soustrait $2x$ des deux côtés)} \\
    3x - 8 &= 7 && \text{(on réduit le membre de gauche)} \\
    3x &= 7 + 8 && \text{(on ajoute $8$ des deux côtés)} \\
    3x &= 15 \\
    x &= \frac{15}{3} = \repbox{5} && \text{(on divise par $3$)}
\end{align*}
\textbf{Vérification} : Pour $x = 5$, Membre de gauche $= 5 \times 5 - 8 = 17$, Membre de droite $= 2 \times 5 + 7 = 17$. L'égalité est vérifiée, la solution unique est bien $x = 5$.
\end{cadreexemple}

\subsection{Équations Produits Nuls}

\begin{cadreprop}[Règle du produit nul]
Un produit de facteurs est nul si et seulement si \textbf{au moins l'un des facteurs est nul} :
\[ A \times B = 0 \ssi A = 0 \quad\text{ou}\quad B = 0 \]
\end{cadreprop}

\begin{cadreexemple}[Résolution d'une équation produit]
Résoudre $(2x - 6)(3x + 12) = 0$ :
\begin{align*}
    2x - 6 = 0 \quad &\text{ou} \quad 3x + 12 = 0 \\
    2x = 6 \quad &\text{ou} \quad 3x = -12 \\
    x = \frac{6}{2} = \repbox{3} \quad &\text{ou} \quad x = \frac{-12}{3} = \repbox{-4}
\end{align*}
L'équation admet donc deux solutions : $x_1 = 3$ et $x_2 = -4$.
\end{cadreexemple}

% ==============================================================================
% 4. AUTOMATISMES & QUESTIONS FLASH
% ==============================================================================
\section{Questions Flash / Automatismes}

\begin{automatisme}[Résolutions Mentales \& Réflexes (10 min)]
\noindent
\begin{tabularx}{\linewidth}{X|X}
\textbf{1.} $x + 7 = 15 \implies x = \pointille[2cm]$ & \textbf{6.} $3x + 4 = 19 \implies x = \pointille[2cm]$ \\
\textbf{2.} $x - 9 = 4 \implies x = \pointille[2cm]$ & \textbf{7.} $5x - 10 = 0 \implies x = \pointille[2cm]$ \\
\textbf{3.} $4x = 28 \implies x = \pointille[2cm]$ & \textbf{8.} $2x = x + 8 \implies x = \pointille[2cm]$ \\
\textbf{4.} $-3x = 18 \implies x = \pointille[2cm]$ & \textbf{9.} $(x - 5)(x + 2) = 0 \implies x = \pointille[1.5cm] \text{ ou } \pointille[1.5cm]$ \\
\textbf{5.} $\frac{x}{5} = 6 \implies x = \pointille[2cm]$ & \textbf{10.} $x = 3$ est-il solution de $2x + 1 = 7$ ? \pointille[1.5cm] \\
\end{tabularx}
\end{automatisme}

% ==============================================================================
% 5. TRAVAUX DIRIGÉS (TD) - 10 EXERCICES CORRIGÉS
% ==============================================================================
\section{Travaux Dirigés (TD) : Exercices d'Entraînement}

\begin{exercice}[2 pts]{\capRealiser}{Équations de base à une étape}
Résoudre les équations suivantes :
\begin{enumerate}
    \item $x + 14 = 31$
    \item $y - 8{,}5 = 12$
    \item $6x = 42$
    \item $-7t = 56$
\end{enumerate}
\lignesreponse[4]
\end{exercice}

\begin{exercice}[3 pts]{\capRealiser}{Équations du type $ax + b = c$}
Résoudre dans $\mathbb{R}$ :
\begin{enumerate}
    \item $3x + 7 = 25$
    \item $5x - 13 = 17$
    \item $-4x + 9 = -11$
    \item $2x + 15 = 4$
\end{enumerate}
\lignesreponse[6]
\end{exercice}

\begin{exercice}[3 pts]{\capRealiser}{Équations avec l'inconnue dans les deux membres}
Résoudre les équations :
\begin{enumerate}
    \item $7x - 4 = 3x + 16$
    \item $8x + 5 = 2x - 19$
    \item $2(3x - 1) = 4x + 10$
\end{enumerate}
\lignesreponse[6]
\end{exercice}

\begin{exercice}[2 pts]{\capRealiser}{Équations Produits Nuls}
Résoudre les équations produit :
\begin{enumerate}
    \item $(x - 7)(2x + 8) = 0$
    \item $(3x - 15)(4x + 2) = 0$
\end{enumerate}
\lignesreponse[4]
\end{exercice}

\begin{exercice}[3 pts]{\capAnalyser}{Problème Plomberie : Facturation de dépannage}
Un plombier applique un forfait déplacement de $45\text{ \euro}$ et un tarif horaire de $38\text{ \euro}$ de l'heure.
Une facture totale s'élève à $159\text{ \euro}$ TTC.
\begin{enumerate}
    \item Écrire une équation traduisant cette situation en notant $h$ le nombre d'heures de main d'œuvre.
    \item Résoudre cette équation et indiquer la durée exacte de l'intervention.
\end{enumerate}
\lignesreponse[4]
\end{exercice}

\begin{exercice}[3 pts]{\capAnalyser}{Problème Mécanique : Équilibre de charges}
Un atelier automobile charge une camionnette. La masse à vide du véhicule est de $1\,450\text{ kg}$. On y charge 12 caisses d'outillage de même masse $m$ et un groupe électrogène de $130\text{ kg}$.
La masse totale du véhicule chargé est mesurée à $2\,300\text{ kg}$.
\begin{enumerate}
    \item Poser l'équation permettant de calculer la masse $m$ d'une caisse.
    \item Résoudre cette équation et conclure.
\end{enumerate}
\lignesreponse[5]
\end{exercice}

% ==============================================================================
% 6. SUJET TYPE DNB PRO
% ==============================================================================
\section{Vers le Brevet (DNB Pro)}

\begin{sujetdnb}[DNB Pro Centres Étrangers][8 pts]{Location de Véhicules Utilitaires}
Pour transporter des matériaux de chantier, une entreprise hésite entre deux agences de location d'utilitaires pour une journée :
\begin{itemize}
    \item \textbf{Agence RapideLoc} : Forfait journalier de $60\text{ \euro}$ plus $0{,}25\text{ \euro}$ par kilomètre parcouru.
    \item \textbf{Agence EcoDrive} : Forfait journalier de $30\text{ \euro}$ plus $0{,}40\text{ \euro}$ par kilomètre parcouru.
\end{itemize}
On note $x$ la distance parcourue en kilomètres au cours de la journée.

\begin{enumerate}
    \item \capSApproprier\ Exprimer en fonction de $x$ le coût de location $C_1(x)$ pour RapideLoc et $C_2(x)$ pour EcoDrive.
    \item \capRealiser\ Calculer le coût pour chaque agence si l'entreprise parcourt une distance de $100\text{ km}$.
    \item \capAnalyser\ Résoudre l'équation $60 + 0{,}25x = 30 + 0{,}40x$.
    \item \capValider\ Interpréter la valeur de $x$ trouvée dans le contexte commercial : pour quelle distance les deux agences facturent-elles exactement le même montant ?
    \item \capCommuniquer\ Conseiller l'entreprise si elle prévoit d'effectuer un trajet de $350\text{ km}$.
\end{enumerate}
\lignesreponse[7]
\end{sujetdnb}

% ==============================================================================
% 7. CORRIGÉS DÉTAILLÉS
% ==============================================================================
\section{Corrigés Détaillés du Chapitre}

\begin{solution}[Corrigé des Exercices de TD]
\begin{itemize}
    \item \textbf{Exercice 1} :
    1. $x = 31 - 14 = \repbox{17}$ \quad 2. $y = 12 + 8{,}5 = \repbox{20{,}5}$ \quad 3. $x = 42/6 = \repbox{7}$ \quad 4. $t = 56/(-7) = \repbox{-8}$.
    \item \textbf{Exercice 2} :
    1. $3x = 25 - 7 = 18 \implies x = \repbox{6}$ \quad 2. $5x = 30 \implies x = \repbox{6}$ \\
    3. $-4x = -20 \implies x = \repbox{5}$ \quad 4. $2x = 4 - 15 = -11 \implies x = \repbox{-5{,}5}$.
    \item \textbf{Exercice 3} :
    1. $7x - 3x = 16 + 4 \implies 4x = 20 \implies x = \repbox{5}$ \\
    2. $8x - 2x = -19 - 5 \implies 6x = -24 \implies x = \repbox{-4}$ \\
    3. $6x - 2 = 4x + 10 \implies 2x = 12 \implies x = \repbox{6}$.
    \item \textbf{Exercice 4} :
    1. $x = 7$ ou $2x = -8 \implies x = \repbox{-4}$. Solutions : $\{ -4 ; 7 \}$. \\
    2. $3x = 15 \implies x = \repbox{5}$ ou $4x = -2 \implies x = \repbox{-0{,}5}$. Solutions : $\{ -0{,}5 ; 5 \}$.
    \item \textbf{Exercice 5 (Plomberie)} :
    1. Équation : $38h + 45 = 159$. \\
    2. $38h = 159 - 45 = 114 \implies h = 114 / 38 = \repbox{3\text{ heures}}$.
    \item \textbf{Exercice 6 (Mécanique)} :
    1. Équation : $1450 + 12m + 130 = 2300 \implies 12m + 1580 = 2300$. \\
    2. $12m = 2300 - 1580 = 720 \implies m = 720 / 12 = \repbox{60\text{ kg}}$.
\end{itemize}
\end{solution}

\begin{solution}[Corrigé du Sujet DNB Pro]
\begin{enumerate}
    \item $C_1(x) = 0{,}25x + 60$ et $C_2(x) = 0{,}40x + 30$.
    \item Pour $x = 100\text{ km}$ :
    $C_1(100) = 0{,}25 \times 100 + 60 = 25 + 60 = 85\text{ \euro}$. \\
    $C_2(100) = 0{,}40 \times 100 + 30 = 40 + 30 = 70\text{ \euro}$. EcoDrive est la moins chère.
    \item $60 - 30 = 0{,}40x - 0{,}25x \implies 30 = 0{,}15x \implies x = \frac{30}{0{,}15} = \repbox{200}$.
    \item Pour une distance de \textbf{200 km}, les deux agences facturent exactement le même tarif ($110\text{ \euro}$).
    \item Pour $x = 350\text{ km} > 200\text{ km}$, c'est l'agence **RapideLoc** qui est la plus économique ($C_1(350) = 147{,}50\text{ \euro}$ contre $C_2(350) = 170\text{ \euro}$).
\end{enumerate}
\end{solution}
